% ============================================================
\section{The Large Discovery Model}
\label{sec:method}
% ============================================================

We are now ready to define the \emph{Large Discovery Model} (LDM). LDM is a
policy for sequential inverse design built from three interacting components:
a generative foundation model, a probabilistic surrogate, and a model-based
acquisition function.

Same as before at iteration $t$, let
\[
\mathcal{D}_t=\{(x_i,r_i)\}_{i=1}^{t-1}
\]
denote the empirical evaluation history, where $x_i$ is a previously evaluated
design and $r_i$ is its observed reward. We use $\mathcal{C}_t$ to denote the
broader \emph{search context} available to the generative model. This context
may include $\mathcal{D}_t$, previously generated candidates, surrogate
predictions, acquisition values, constraint feedback, refinement trajectories,
and other forms of search memory. Thus, $\mathcal{D}_t$ contains the empirical
evidence used to fit the surrogate, whereas $\mathcal{C}_t$ contains the
information used to condition candidate generation.

\paragraph{(i) A generative foundation model.}
We denote by
\begin{equation}
    p_{\theta,\alpha}(x\mid\mathcal{C}_t)
\end{equation}
the proposal distribution over plausible designs conditioned on the current
search context. Here, $\theta$ denotes the parameters of the underlying
generative model, while $\alpha$ denotes its inference configuration, such as
the prompting strategy, sampling temperature, reasoning procedure, refinement
scheme, or inference-time compute budget. We use $p_{\theta,\alpha}$ throughout
the paper to make explicit that the effective proposal distribution depends
not only on the model parameters but also on how the model is used at
inference time.

The generative model supplies a structured prior over the design space
$\mathcal{X}$. It enables the search to generate candidates that reflect
domain knowledge, semantic plausibility, and structural constraints without
requiring $\mathcal{X}$ to be explicitly enumerated. By changing its context
or inference procedure, the model can also expand and reshape the set of
candidates reachable by the current search.

\paragraph{(ii) A probabilistic surrogate.}
As introduced in Section~\ref{sec:ldm}, the posterior
\[
    p(R\mid\mathcal{D}_t)
\]
represents the current belief about the unknown reward function, conditioned
on the empirical observations collected so far. For a candidate $x$, the
predictive mean $\mu_t(x)$ estimates its expected reward, while the predictive
uncertainty $\sigma_t(x)$ quantifies the surrogate's epistemic uncertainty
about that estimate. High epistemic uncertainty indicates that the prediction
is weakly supported by the available evidence and may be reduced through
additional informative evaluations. It is therefore distinct from
irreducible randomness or measurement noise in the evaluation process.

The surrogate provides the empirical grounding that the generative model
alone generally lacks. Although a foundation model may encode substantial
domain knowledge, its likelihoods and self-assessments do not necessarily
provide calibrated estimates of a task-specific quantitative reward,
particularly for novel candidates outside the observed data distribution.

\paragraph{(iii) A model-based acquisition function.}
From the surrogate's predictive distribution, we construct an acquisition
function
\[
    \acq_t(x)
\]
that assigns each candidate a scalar value to the sequential search process.
Depending on its form, the acquisition function may favour candidates with
high predicted reward, candidates with high epistemic uncertainty, or a
principled combination of the two. It therefore mediates the trade-off between
exploitation and exploration.

The acquisition value should not be interpreted simply as the predicted
reward of a candidate. Rather, it measures the potential value of allocating
further computation or empirical evaluation to that candidate. It can be used
both to decide which candidates warrant costly external evaluation and to
direct inference-time operations such as sampling, ranking, refinement, and
branch expansion.

These three components play complementary roles. The generative model proposes
structured candidates and can move the search frontier towards designs not
previously considered. The surrogate grounds the search in empirical
observations and quantifies epistemic uncertainty. The acquisition function
uses this predictive belief to determine which reachable candidates are most
valuable to develop or evaluate. LDM thereby uses empirically calibrated
feedback to direct the inference-time search of the generative model.

%\subsection{Acquisition-Tilted Search}

We seek a search distribution that attains high expected acquisition value
while remaining sufficiently close to the structured prior supplied by the
generative model. Assume that $\mathcal{X}$ is a measurable space, with its
$\sigma$-algebra omitted for simplicity. At iteration $t$, we define
\begin{equation}
\label{eq:variational}
    \pi_t
    =
    \argmax_{q\in\Delta(\mathcal{X})}
    \left\{
        \mathbb{E}_{x\sim q}\!\left[\acq_t(x)\right]
        -
        \frac{1}{\eta}
        \KL\!\left(
            q
            \,\middle\|\,
            p_{\theta,\alpha}(\cdot\mid\mathcal{C}_t)
        \right)
    \right\},
\end{equation}
where $\Delta(\mathcal{X})$ denotes the set of probability measures over
$\mathcal{X}$, $\KL$ is the Kullback--Leibler divergence, and $\eta>0$ controls
the strength of acquisition-based tilting. The expected-acquisition term moves
probability mass towards candidates that are valuable under the empirically
grounded surrogate. The KL term prevents the search policy from departing
arbitrarily far from the generative model's prior over plausible and
well-formed designs. The parameter $\eta$ controls this trade-off: larger
values place greater emphasis on acquisition value, whereas smaller values
keep the search closer to the generative prior.

The inference configuration $\alpha$ affects $\pi_t$ through
$p_{\theta,\alpha}(\cdot\mid\mathcal{C}_t)$ and therefore changes the reference
distribution against which the KL divergence is measured. Increasing
inference-time computation may reshape this proposal distribution through
additional sampling, reasoning, refinement, or structured search. The
acquisition function then determines how that computation is allocated
according to predicted empirical value and epistemic uncertainty.

Equation~\eqref{eq:variational} is deliberately model-agnostic: it does not
depend on a particular generative architecture, probabilistic surrogate, or
acquisition function. Its variational form follows the Gibbs variational principle and is closely
related to Gibbs posteriors~\citep{donsker1975variational,bissiri2016general},
KL-regularised policy search~\citep{peters2010relative},
maximum-entropy reinforcement learning~\citep{haarnoja2018soft}, and control as
probabilistic inference~\citep{levine2018reinforcement}. Accordingly, our contribution is not the generic
KL-regularised objective in isolation. The distinctive role of
Equation~\eqref{eq:variational} in LDM is to connect a structured generative
prior to an acquisition value derived from a continually updated empirical
surrogate. This connection allows external observations to direct both
inference-time computation and sequential evaluation over structured,
potentially open-ended design spaces.

The resulting policy is recurrent (see Fig~\ref{fig:loop}). Candidates generated under the current
search context are scored and refined using the surrogate-derived acquisition
signal; selected candidates are then evaluated by an external empirical
source; and the resulting observations update $\mathcal{D}_t$, the surrogate,
and the next search context $\mathcal{C}_{t+1}$. LDM therefore does not merely
reweight a fixed set of generated samples. It repeatedly couples generation,
empirical evaluation, belief updating, and search-frontier expansion.

In the following sections, we first show that
Equation~\eqref{eq:variational} admits a closed-form acquisition-tilted policy.
We then describe how this ideal policy can be approximated through
inference-time sampling, selection, and iterative refinement, with particular
attention to test-time compute scaling. We subsequently instantiate the
probabilistic surrogate as a Gaussian process and construct $\acq_t$ from its
posterior predictive distribution. Finally, we describe how the same
acquisition signal can support parameter adaptation: in addition to directing
inference through $\alpha$, it may provide supervision for updating the
generative parameters $\theta$ through fine-tuning.


\subsection{An Acquisition-Tilted Search Distribution}
\label{sec:core}

With the LLM prior $p_{\theta,\alpha}$ and the acquisition function $\acq_t$ defined, our desired search policy $\pi_t$ follows directly from the variational objective in Eq.~\eqref{eq:variational}. This variational problem is a canonical instance of the Gibbs variational principle (equivalently, the Donsker–Varadhan representation of the KL divergence)~\citep{donsker1975variational,bissiri2016general}, whose unique maximiser is the exponentially tilted distribution. This form of KL-regularised reward maximisation is the same mathematical foundation underlying maximum-entropy reinforcement learning, control-as-inference \citep{levine2018reinforcement}, and the closed-form policy objective in KL-regularised RLHF \citep{rafailov2023direct}.

\begin{proposition}[Optimal search distribution]
\label{prop:gibbs}
The unique maximiser of \eqref{eq:variational} is
\begin{equation}
\label{eq:core}
\pi_t(x)=\frac{1}{Z_t}\,p_{\theta,\alpha}(x\mid\mathcal{C}_t)\,\exp\!\big\{\eta\,\acq_t(x)\big\},
\qquad Z_t=\int_{x\in\mathcal{X}}p_{\theta,\alpha}(x\mid\mathcal{C}_t)\,\exp\!\big\{\eta\,\acq_t(x)\big\}\,\mathrm{d}x.
\end{equation}
\end{proposition}
\begin{proof}
Please refer to Appendix~\ref{prop_proof}.
\end{proof}

Eq.~\eqref{eq:core} shows that the optimal policy $\pi_t$ takes an acquisition-tilted form, where the base measure is reweighted by the exponent of the acquisition function. Interestingly, this policy draws an analogy to the infinite-armed bandit framework \citep{berry1997bandit,wang2008infinitely}: as scientific designs (arms) cannot be listed in advance, they are revealed from a reservoir distribution, and the policy tilts that reservoir toward arms it judges valuable. In this reading, $p_{\theta,\alpha}(\cdot\mid\mathcal{C}_t)$ plays the role of the reservoir, $\eta\,\acq_t(x)$ acts as the strength-of-reward tilt, and $Z_t$ is the usual log-sum-exp normaliser.

A clear interpretation of the inference configuration $\alpha$ and the KL-regularisation coefficient $\eta$ can be drawn through standard analysis of infinite-armed bandits. We restrict attention to the common case in which the reservoir admits the temperature form
\begin{equation}
  \label{eq:temp}
  p_{\theta,\alpha}(x\mid\mathcal{C}_t) \propto p_\theta(x\mid\mathcal{C}_t)^\alpha,
\end{equation}
where $p_\theta$ denotes the unconditional base proposal. This is the exponential family of reservoirs commonly used in the bandit literature, and it makes $\alpha$ transparent as the strength of the prior: large values of $\alpha$ concentrate probability mass on modes of $p_\theta$, while $\alpha\to 0$ spreads mass toward a uniform measure on $\mathcal{X}$. Substituting this form into \eqref{eq:core} and taking the mode gives the bandit-style argmax
\begin{equation}
  \label{eq:argmax}
  x_{t+1}=\argmax_{x\in\mathcal{X}}\ \Big[\ \alpha\,\log p_\theta(x\mid\mathcal{C}_t)+\eta\,\acq_t(x)\ \Big],
\end{equation}
which picks the arm that maximises a weighted combination of the prior log-mass $\log p_\theta$ and the reward tilt $\acq_t$, i.e., the bandit analogue of selecting the most promising revealed arm.

Examining the two limiting cases of the temperature form $p_{\theta,\alpha}(x\mid\mathcal{C}_t)\propto p_\theta(x\mid\mathcal{C}_t)^\alpha$ situates LDM relative to related paradigms. As $\eta\to 0$ the tilt vanishes and $\pi_t$ collapses to the raw reservoir $p_{\theta,\alpha}(\cdot\mid\mathcal{C}_t)$, so search is driven entirely by the generative model with no data feedback. As $\alpha\to0$ the reservoir spreads toward a uniform measure on $\mathcal{X}$ and the policy becomes $\pi_t\propto\exp\{\eta\,\acq_t(\cdot)\}$, recovering classical Bayesian optimisation once $\eta$ is also taken large. LDM is the regime between these two limit points, where both the structured prior and the uncertainty-aware acquisition contribute to the policy. We note, however, that some LLM inference backends do not admit the temperature form $p_\theta(x\mid\mathcal{C}_t)^\alpha$, and the preceding analysis should be read as the canonical case under which the role of $\alpha$ is most cleanly interpreted.

With this construction, LDM defines a discovery policy $\pi_t$ from which new designs may be drawn using either a stochastic or a deterministic approach, as illustrated by Eq.~\eqref{eq:argmax}. In practice, the normalising constant $Z_t$ is never computed explicitly, and we instead rely on inference-time search to approximate draws from the tilted distribution.

\subsection{Surrogate Modelling of the Reward Function}
\label{sec:surrogate}

To efficiently navigate the design space, we require a probabilistic model over function space that treats the reward function $R$ as a random function and models the posterior belief $R \mid \mathcal{D}_t$ conditioned on our cumulative evaluations. Our LDM formulation is largely inspired by the uncertainty-aware decision-making paradigm of Bayesian optimisation \citep{garnettBayesianOptimization2023}.

We model this objective using a Gaussian Process (GP) \citep{rasmussen2006gaussian}. GPs serve as highly effective surrogates in this setting because they are exceptionally sample-efficient, admit exact Bayesian updates, and provide principled non-parametric uncertainty estimates across unexplored regions of the design space. Formally, we define the prior distribution over the reward function as:
\begin{equation}
\label{eq:gp_prior_main}
R(x) \sim \mathcal{GP}\bigl(m(x),\, k(x, x')\bigr),
\end{equation}
where $m: \mathcal{X} \to \mathbb{R}$ is the prior mean function and $k: \mathcal{X} \times \mathcal{X} \to \mathbb{R}$ is a positive-definite kernel function capturing spatial correlation. 

By conditioning this prior on the historical data $\mathcal{D}_t = \{(x_i, r_i)\}_{i=1}^{t}$, we obtain a Gaussian posterior distribution characterised by a predictive mean $\mu_t(x)$ and a predictive variance $\sigma_t^2(x)$ for any candidate point $x \in \mathcal{X}$. These two statistical moments allow us to construct acquisition functions that elegantly balance exploitation (seeking high predicted rewards) and exploration (targeting regions of high uncertainty). For instance, the popular Upper Confidence Bound (UCB) acquisition function is formulated as:
\begin{equation}
\label{eq:ucb_main}
\alpha^{\mathrm{UCB}}_t(x) = \mu_t(x) + \sqrt{\beta_t}\,\sigma_t(x),
\end{equation}
where $\beta_t > 0$ is a parameter scaling the exploration incentive. The complete technical details of the posterior inference, along with the formulation of other common acquisition functions, are provided in Appendix~\ref{sec:llm-bo}.

Because UCB decomposes the decision value into a high-mean and a high-uncertainty term, the tilted policy $\pi_t(x)\propto p_{\theta,\alpha}(x\mid\C_t)\exp\{\eta\,\acq_t(x)\}$ of Eq.~\eqref{eq:core} instantiates the discovery--exploration--exploitation triangle of Figure~\ref{fig:knowledge-regimes} directly: the LLM reservoir term realises discovery, the high-$\sigma_t$ tail of $\acq_t$ realises exploration, and the high-$\mu_t$ mode realises exploitation. We use UCB as the running example precisely because the trade-off is explicit in its closed form; in richer acquisitions the same triangle is realised only implicitly through the surrogate.

\subsection{Acquisition-Tilted Inference-Time Search}
\label{sec:inference-time-search}

Exact evaluation or sampling of $\pi_t$ is computationally intractable over the high-dimensional, combinatorial design spaces typical of scientific discovery. We therefore approximate the tilted distribution via inference-time search: the LLM proposes a finite pool of candidates, which are then reweighted by the acquisition function to align with the target distribution. This framework naturally supports test-time compute scaling, mirroring the broader paradigm in general LLM reasoning but guided by a calibrated, task-grounded value signal rather than model-internal confidence~\citep{wang2024openr}.

We consider two canonical approximation schemes. The deterministic best-of-$N$ rule targets the mode of the tilted distribution: we draw $N$ independent candidates from the LLM prior,
\[
\{\hat{x}_{t,i}\}_{i=1}^N \sim p_{\theta, \alpha}(\cdot \mid \mathcal{C}_t),
\]
and select the candidate with the highest acquisition score:
\[
x_{t+1} = \argmax_{i \in \{1,\dots,N\}} \acq_t(\hat{x}_{t,i}).
\]
The proposal $p_{\theta,\alpha}$ may comprise a collection of hyperparameters, e.g.,\ temperature and top-$p$ for an autoregressive decoder, or a noise schedule for a diffusion sampler, depending on the backend. The single explicit tunable below is $\eta$; the rest of $p_{\theta,\alpha}$ is held fixed.

The stochastic scheme approximates sampling from the full tilted distribution via softmax reweighting, where sampling probability is proportional to the exponentiated acquisition value:
\begin{equation}
\label{eq:softmax_reweight}
\Pr(x_t = \hat{x}_{t,i})
=
\frac{\exp\bigl\{\eta \, \acq_t(\hat{x}_{t,i})\bigr\}}
{\sum_{j=1}^N \exp\bigl\{\eta \, \acq_t(\hat{x}_{t,j})\bigr\}}.
\end{equation}
For batch selection with $b>1$, these generalise naturally to top-$b$ selection for the deterministic case and Gumbel-top-$b$ sampling for the stochastic case.

Candidate pool size $N$ plays the critical role of \emph{test-time compute scaling}, a powerful technique for LLM reasoning \citep{wangTutorialLLMReasoning2025}, and LDM adapts this core advantage to discovery tasks. Larger $N$ improves the fidelity of the approximation to the optimal tilted distribution and lifts empirical performance, at the cost of additional LLM inference and surrogate evaluation. This aligns with the established finding that test-time compute can substitute for model scale in LLM reasoning, with the key distinction that LDM uses an externally calibrated acquisition function as guidance, rather than the model’s own self-evaluation or linguistic plausibility.

The full-candidate reweighting scheme is simple, modality-agnostic, and forms the basis of our experiments. For sequentially generated design spaces such as sequences and programs, the principle can be extended to intermediate generation steps via a prefix value function $V(x_{t,\le i}) = \mathbb{E}_{x_{t,>i}}[\acq_t(x_{t,\le i},x_{t,>i})]$, enabling more advanced methods such as beam search and Monte Carlo Tree Search~\citep{fengAlphazerolikeTreeSearchCan2024}. These multi-step variants offer further search coverage at higher computational cost, and we leave their full empirical evaluation as future work.

\begin{figure}[t]
\centering
\begingroup
\definecolor{llmred}{RGB}{249,220,220}%
\definecolor{llmredline}{RGB}{201,74,74}%
\definecolor{gpblue}{RGB}{214,231,247}%
\definecolor{gpblueline}{RGB}{58,110,165}%
\definecolor{graybox}{RGB}{238,238,238}%
\definecolor{grayline}{RGB}{120,120,120}%
\definecolor{fbgray}{RGB}{90,90,90}%
\resizebox{\textwidth}{!}{%
\begin{tikzpicture}[
    >={Latex[length=2.4mm]},
    font=\normalsize,
    llm/.style={draw=llmredline,line width=1pt,fill=llmred,rounded corners=5pt,
                align=center,inner sep=7pt,minimum height=20mm,text width=50mm},
    gp/.style={draw=gpblueline,line width=1pt,fill=gpblue,rounded corners=5pt,
                align=center,inner sep=7pt,minimum height=20mm,text width=50mm},
    proc/.style={draw=grayline,line width=0.9pt,fill=graybox,rounded corners=4pt,
                align=center,inner sep=6pt},
    lbl/.style={font=\small,align=center},
    fb/.style={font=\small\itshape,align=center,text=fbgray},
    flow/.style={->,line width=0.9pt},
    fbflow/.style={->,line width=1.0pt,fbgray},
]
 
% Top boxes
\node[llm] (llm) at (-2.9,3.7)
  {\textbf{LLM reservoir}\\[3pt]$p_{\theta,\alpha}(\cdot\mid\C_t)$};
\node[gp]  (gp)  at ( 2.9,3.7)
  {\textbf{GP surrogate: reward model}\\[3pt]$\mu_t,\sigma_t$, acquisition $\acq_t$};
 
% Tilted-search box
\node[proc,text width=78mm] (pi) at (0,0.65)
  {\textbf{acquisition-tilted search}\\[3pt]
   $\pi_t(x)\;\propto\;p_{\theta,\alpha}(x\mid\C_t)\,\exp\!\big\{\eta\,\acq_t(x)\big\}$};
 
% Candidate batch
\node[proc,text width=70mm] (cand) at (0,-1.65)
  {\textbf{draw candidate batch }$B_t=\{x_{t,1},\dots,x_{t,b}\}$};
 
% Evaluate batch
\node[proc,text width=68mm] (eval) at (0,-3.95)
  {\textbf{evaluate }$R(x)$\textbf{ for all }$x\in B_t$};
 
% Arrows: top boxes into pi
\draw[flow] (llm.south) -- (pi.north -| llm.south)
   node[lbl,midway,left=3pt] {base measure};
\draw[flow] (gp.south) -- (pi.north -| gp.south)
   node[lbl,midway,right=3pt] {acquisition function};
 
% pi -> cand -> eval
\draw[flow] (pi) -- (cand)
   node[lbl,midway,right=1pt] {inference-time search};
\draw[flow] (cand) -- (eval)
   node[lbl,midway,right=1pt] {query};
 
% Feedback loops (both labelled "feedback")
% eval -> llm  (update history)
\draw[fbflow] (eval.west) .. controls +(-3.6,0) and +(-1.8,0) .. (llm.west)
   node[fb,pos=0.28,fill=white,inner sep=2pt] {feedback}
   node[lbl,pos=0.62,fill=white,inner sep=2pt] {update history $\C_t$};
% eval -> gp  (update data)
\draw[fbflow] (eval.east) .. controls +(3.6,0) and +(1.8,0) .. (gp.east)
   node[fb,pos=0.28,fill=white,inner sep=2pt] {feedback}
   node[lbl,pos=0.62,fill=white,inner sep=2pt] {update data $\D_t$};
 
\end{tikzpicture}%
}
\endgroup
\caption{The Large Discovery Model loop. The LLM supplies the base measure
$p_{\theta,\alpha}(\cdot\mid\C_t)$ and the GP surrogate supplies the acquisition function
$\acq_t$; together they define the acquisition-tilted search distribution $\pi_t$. At each
round, inference-time search draws a candidate batch $B_t=\{x_{t,1},\dots,x_{t,b}\}$,
evaluates $R(x)$ for every $x\in B_t$, and feeds the observations back to update data $\D_t$
and history $\C_t$.}
\label{fig:loop}
\end{figure}


% ============================================================
\subsection{Fine-Tuning to Reduce Test-Time Search}
\label{sec:ldm-finetune}
% ============================================================

The LDM policy in Equation~\eqref{eq:core} combines the structured generative
prior (reservoir)
$p_{\theta,\alpha}(\cdot\mid\mathcal{C}_t)$
with the acquisition function $\acq_t$. A high-budget implementation of LDM
test-time search, denoted by LDM-TTS, approximates this policy by generating a
large pool of candidates and using the acquisition function to rank, select,
and refine the most valuable ones. Although effective, repeating this
procedure at every search round can require substantial inference-time
computation. We therefore propose to distil the behaviour of high-budget
LDM-TTS into a more efficient student proposer through fine-tuning.

At iteration $t$, the acquisition value of a candidate $x$ may be written
abstractly as
\begin{equation}
\label{eq:sft-acquisition}
    \acq_t(x)
    =
    \operatorname{Acq}\!\left(
        \mu_t(x),
        \sigma_t(x),
        g_t(x)
    \right),
\end{equation}
where $\mu_t(x)$ and $\sigma_t(x)$ are the surrogate's posterior predictive
mean and epistemic uncertainty, respectively. The term $g_t(x)$ collects
additional task-dependent considerations, such as constraint satisfaction,
novelty, evaluation cost, diversity, or expected improvement of the current
Pareto frontier.

Suppose that high-budget LDM-TTS generates a candidate pool
\[
    \widehat{\mathcal{X}}_t
    =
    \{\hat{x}_{t,i}\}_{i=1}^{N}
\]
under the search context $\mathcal{C}_t$. We define an acquisition-weighted
empirical distribution over this pool by
\begin{equation}
\label{eq:empirical-tilt}
    \widehat{\pi}_t(\hat{x}_{t,i}\mid\mathcal{C}_t)
    =
    \frac{
        \exp\!\left\{\eta\,\acq_t(\hat{x}_{t,i})\right\}
    }{
        \sum_{j=1}^{N}
        \exp\!\left\{\eta\,\acq_t(\hat{x}_{t,j})\right\}
    },
\end{equation}
where $\eta>0$ controls how strongly the distribution concentrates on
candidates with high acquisition values. This finite-pool distribution
approximates the ideal acquisition-tilted policy in
Equation~\eqref{eq:core}.

Let $p_{\phi}(x\mid\mathcal{C}_t)$ denote the student proposer, initialised
from the underlying generative model. We fine-tune it by minimising the
acquisition-weighted negative log-likelihood
\begin{equation}
\label{eq:acquisition-distillation}
    \mathcal{L}_{\mathrm{distill}}(\phi)
    =
    -
    \sum_t
    \sum_{i=1}^{N}
    \widehat{\pi}_t(\hat{x}_{t,i}\mid\mathcal{C}_t)
    \log p_{\phi}(\hat{x}_{t,i}\mid\mathcal{C}_t).
\end{equation}
Equivalently, this objective projects the finite-pool teacher policy
$\widehat{\pi}_t$ onto the family of distributions represented by the student
model. The resulting proposer assigns greater probability to candidates that
high-budget LDM-TTS would regard as valuable, allowing some of the benefits of
test-time search to be transferred into the model parameters.

Importantly, the training signal is richer than the single candidate
eventually selected for costly empirical evaluation. The unselected candidates
record which alternatives the generative model considered and how the
surrogate ranked their exploitation, exploration, and discovery value under
the same search context. Fine-tuning on this information amortises part of the
high-budget search process: the student learns to propose stronger candidates
directly and can therefore achieve comparable search quality with a smaller
candidate pool or fewer refinement steps at test time.

Because the surrogate and acquisition function change as new empirical
observations are collected, the distilled policy is associated with the
search states on which it was trained. It may therefore be updated
periodically as the empirical dataset and search frontier evolve, rather than
being treated as a permanently fixed replacement for online LDM search.
% Insert the existing Figure~\ref{fig:training-pipeline} TikZ block here.
% Keep its label and use the following concise caption.
%
% \caption{
% \textbf{Amortizing LDM test-time search.}
% High-budget LDM-TTS samples a candidate pool from the LLM reservoir and scores
% it with the current surrogate acquisition function. High-value search traces,
% optionally paired with surrogate-grounded rationales, are distilled through
% supervised fine-tuning into a low-budget discovery proposer.
% }
% \label{fig:training-pipeline}

\paragraph{Reasoning augmentation.}
Acquisition values rank candidates but do not explicitly state why a candidate
is useful for research progress. We therefore optionally augment selected
teacher traces with a short rationale
\begin{equation}
z_{t,i}
=
\operatorname{Explain}\Big(
\C_t,\hat{x}_{t,i},
\mu_t(\hat{x}_{t,i}),
\sigma_t(\hat{x}_{t,i}),
g_t(\hat{x}_{t,i}),
\acq_t(\hat{x}_{t,i})
\Big).
\label{eq:reasoning-augmentation}
\end{equation}
The rationale is grounded in the surrogate value decomposition. It can explain
whether a candidate exploits high predicted utility, probes an uncertain
region, satisfies a constraint, preserves diversity, or tests a new candidate
family after search stagnation. Thus, unlike reasoning traces for closed-world
problems, discovery rationales explain why a scarce experiment is worth
running.

We construct an SFT dataset of tuples
$(\C_t,z_{t,i},\hat{x}_{t,i})$, retaining or sampling candidate traces
according to $\widehat{\pi}^{(N)}_t$. Let $\bar{w}_{t,i}$ be normalized weights
derived from Eq.~\eqref{eq:empirical-tilt}. A student model $p_\phi$ is trained
with the weighted autoregressive objective
\begin{equation}
\mathcal{L}_{\mathrm{SFT}}(\phi)
=
-
\sum_t\sum_{i=1}^{N}
\bar{w}_{t,i}
\left[
\log p_\phi(z_{t,i}\mid\C_t)
+
\log p_\phi(\hat{x}_{t,i}\mid\C_t,z_{t,i})
\right].
\label{eq:sft-loss}
\end{equation}
Equivalently, trajectories may be sampled from
$\widehat{\pi}^{(N)}_t$ and optimised with an unweighted likelihood objective.
The rationale is an auxiliary training target and can be omitted at
deployment.

Fine-tuning shifts the student reservoir toward candidate families that
high-budget LDM-TTS repeatedly identifies as acquisition-relevant. In the
regret decomposition of Section~\ref{sec:theory-main}, this can reduce the
discovery gap by improving reservoir coverage and reduce the LDM sampling
shortfall by increasing the near-acquisition reservoir mass
$\kappa_t(\zeta_t)$. Fine-tuning does not replace the online surrogate:
new observations must still update $\mu_t$, $\sigma_t$, and $\acq_t$. Instead,
it reduces the candidate-pool size and inner-search budget needed before the
surrogate can identify valuable experimental actions.